\section{Discussion and Conclusion}
\label{sec:discussion}

\paragraph{Methodological lessons for the steering literature.}
Three of our findings generalize beyond probe prompting and beyond Gemma-2-2B-it.
\emph{(i) Match random controls structurally.} Without matching feature count, per-layer histogram, and concept-exclusion, the labeled--random gap is over- or under-stated; we observed that random controls produce \emph{more} suppression than labeled in 3/5 domains, so any metric that treats suppression as ``circuit response'' will mistake generic disruption for specific targeting. \emph{(ii) Default $M{=}20$ does not transfer.} The single-feature paradigm's $M{=}20$ default is over-steering in the multi-feature CLT setting; the winning $M$ is bimodal at $\sim 2.4$/$\sim 6.9$. We recommend a coarse $M$ sweep as a default reporting practice. \emph{(iii) Decompose by semantic field role.} The less-is-more effect (intermediate$+$answer beats all-3-fields by 14--33 pp in 4 domains) suggests that input-field features encode what the model \emph{reads}, not what it should \emph{produce}, and disturbing them activates competing circuits.

\paragraph{What the labels mean.}
The supernode labels are behavioral abstractions (\cref{def:opuseful}); they pass an operational test against matched controls but do not constitute mechanistic identifications of latent variables (Level~3 of \citealp{geiger2025causal}). The four functional roles (\textsc{Sem-Dict/Sem-Conc/Rel/Say-X}) may correspond to natural computational types or to convenient partitions of a continuous behavioral space; we cannot distinguish these with the present evidence.

\paragraph{Future work.}
We pre-register four directions: (a) attention-aware variants that re-attribute under free attention; (b) cross-architecture replication on Gemma-2-9B and Llama-3 with available transcoders; (c) multilingual extension with language-specific functional vocabularies; (d) a within-domain pair-level scaffold predictor that exceeds the null observed in USA, possibly by combining scaffold with CLT error-node density and $K$. We commit to releasing all artifacts (\cref{sec:repro}) so that the matched-control harness can be re-run by other groups against alternative feature-grouping methods.

\paragraph{Conclusion.}
Probe prompting demonstrates that circuit-level interpretability of attribution graphs can be partially automated using a transparent rule-based pipeline whose outputs are validated by intervention against matched random controls. The two-condition definition of an operationally useful label (Specificity $+$ Matched-Control) gives the harness a sharp falsification criterion; the harness's substantive findings (less-is-more, over-steering, scaffold) generalize beyond our pipeline and inform best practice in the steering literature.
